% Generated draft: requires mathematical and source review.
\documentclass[11pt]{article}
\usepackage[margin=25mm]{geometry}
\usepackage{fontspec}
\setmainfont{Latin Modern Roman}
\usepackage{amsmath,amssymb,mathtools,unicode-math}
\setmathfont{Latin Modern Math}
\usepackage{enumitem,needspace,xurl,hyperref}
\setlength{\parindent}{0pt}
\setlength{\parskip}{5pt}
\setlength{\emergencystretch}{3em}
\newcommand{\N}{\mathbb{N}}
\newcommand{\Z}{\mathbb{Z}}
\newcommand{\Q}{\mathbb{Q}}
\newcommand{\R}{\mathbb{R}}
\newcommand{\C}{\mathbb{C}}
\newcommand{\F}{\mathbb{F}}
\newcommand{\A}{\mathbb{A}}
\newcommand{\PP}{\mathbb{P}}
\newcommand{\E}{\mathbb{E}}
\newcommand{\eps}{\varepsilon}
\newcommand{\dd}{\,\mathrm{d}}
\begin{document}
\section*{UnsolvedMath ID: 30005095}
\textbf{Source problem code:} catalogue entry
\subsection{Definitions and mathematical statement}
The supplied record concerns an abelian scheme \(\pi:A\to S\), where the original excerpt specifies that \(S\) is an irreducible \(\Q\)-variety, and an irreducible closed subvariety \(X\subseteq A\) satisfying \(\pi(X)=S\). The excerpt does not define the canonical height or a horizontal torsion coset. Accordingly, the following explicit formulation is conditional on the conventions below, which require confirmation against the complete mathematical source.

Choose a symmetric relatively ample line bundle \(\mathcal L\) on \(A\). For an algebraic point \(x\in A(\overline{\Q})\), put \(s=\pi(x)\), let \(A_s\) be the fiber over \(s\), and define
\[
\widehat h_{\mathcal L}(x)=\widehat h_{\mathcal L|_{A_s}}(x),
\]
where the right side is the N\'eron--Tate canonical height on the abelian variety \(A_s\). For each real \(\varepsilon>0\), set
\[
X_{\leq\varepsilon}=\{x\in X(\overline{\Q}):\widehat h_{\mathcal L}(x)\leq\varepsilon\}.
\]
Having a Zariski-dense set of points of arbitrarily small height means that \(X_{\leq\varepsilon}\) is Zariski dense in \(X_{\overline{\Q}}\) for every \(\varepsilon>0\).

For a provisional geometric interpretation of the other alternative, call an irreducible closed subvariety \(T\subseteq A\) a horizontal torsion coset if \(\pi(T)=S\) and every irreducible component of its geometric generic fiber is a translate \(B+t\), where \(B\) is an abelian subvariety of \(A_{\overline\eta}\), \(t\) is a torsion point of \(A_{\overline\eta}\), and \(\overline\eta\) is a geometric generic point of \(S\). This is an explicit candidate convention, not a definition recovered from the supplied excerpt.

Under these conventions, the proposed relative Bogomolov statement is that, for every such \(\pi\) and \(X\),
\[
\left(\exists\varepsilon>0:\ 
\overline{X_{\leq\varepsilon}}^{\,\mathrm{Zar}}
\subsetneq X_{\overline{\Q}}\right)
\quad\text{or}\quad
\left(\exists T\subseteq A:\ T\text{ is a horizontal torsion coset},\ 
X\subseteq T,\ 
\dim T-\dim X\leq\dim S\right).
\]
Here \(\dim T-\dim X\) is the codimension of \(X\) in the containing coset \(T\). The alternatives are not asserted to be mutually exclusive. In particular, the second alternative does not require \(X\) itself to be a horizontal torsion coset. The polarization convention, geometric density convention, and precise intended definition of horizontal torsion coset remain unresolved source qualifications.
\subsection{Short English statement}
For an irreducible subvariety dominating the base of an abelian scheme, the proposed dichotomy says that some positive height threshold gives a non-dense set of points, or the subvariety lies in a horizontal torsion coset with codimension at most the dimension of the base. The explicit formulation is conditional on height and torsion-coset conventions omitted from the supplied excerpt.
\subsection{Sources}
\begin{itemize}
\item Supplied local UnsolvedMath registry record, permanent ID \(30005095\), title \emph{Relative Bogomolov Conjecture for Abelian Schemes}, collection \emph{Oberwolfach Reports Open Problems}, exact source problem code \texttt{catalogue entry}. The record gives proposed year \(2022\), no named proposer, status \texttt{partially\_solved}, and statement status \texttt{corrected\_verified}. Its assertion that the statement was checked against the official report is supplied provenance, not an independent verification performed here.
\item Supplied report citation: \emph{Diophantische Approximationen (2023)}, Oberwolfach Reports, pp.\ 1091--1164, DOI \texttt{10.4171/owr/2022/21}; \url{https://doi.org/10.4171/owr/2022/21}. The supplied excerpt identifies Lars K\"uhne's contribution, \emph{The relative Bogomolov conjecture for fibered products of elliptic families}.
\item Supplied literature reference: K\"uhne, relative Bogomolov for elliptic curves; \url{https://doi.org/10.1515/crelle-2022-0082}.
\item Supplied literature reference: Dimitrov--Gao--Habegger, relative Bogomolov; \url{https://doi.org/10.1016/j.jnt.2021.03.028}.
\item Literature assessment checked at \texttt{2026-08-22}, copied exactly from the record: the relative Bogomolov conjecture is described as proved for arbitrary subvarieties in fiber products of elliptic families, while the general abelian-scheme formulation is described as remaining open. This is supplied dated curation derived from the OWR material, not independently verified mathematical or current literature status.
\item Historical attribution in the supplied excerpt: an earlier, unspecified absolute formulation is credited to Ullmo [11] and Zhang [12], with a later proof by David and Philippon [3]. The excerpt describes the relative conjecture as widely open at the time and reports a workshop announcement of a proof of relative Manin--Mumford by Gao and Habegger [8]. The bibliographic entries behind these reference labels were not supplied, and these historical remarks have not been independently verified.
\end{itemize}
\end{document}
